\section{Discussion}
\label{sec:discussion}

\subsection{Relation to the companion arc}
\label{sec:companions}

This paper is the temporal capstone of a four-part arc built on one open engine. P1
\citep{baweja2026datum} classifies the \emph{spatial} datum defect of a lunar network and
the observation geometry (libration, an absolute tie) that resolves it; its Helmert-seven
datum is the spatial block $\Ssp$ here, imported unchanged, including the SI result that a
geometrically diverse network observes scale. P2 \citep{baweja2026interop} quantifies the
irreducible cross-provider interoperability floor---the inter-ephemeris orientation and
position disagreement between independent authoritative realizations. P3
\citep{baweja2026commonmode} shows that the rigid-body frame freedom is exactly the
error class a snapshot RAIM/ARAIM residual test cannot see (the datum defect is a subset of
the integrity blind subspace). The present paper (P4) adds the missing dimension---%
\emph{time}---and the \emph{coupling}: it extends the datum from the spatial seven to the
coupled nine, classifies the joint null space in basis-invariant form, and shows the joint
gauge is not the direct sum of the spatial and temporal gauges. The four share the same
firewall (exactly one externally validated row per paper) and the same open substrate.

\subsection{Implications for LunaNet, Moonlight/LCNS, and the IAU lunar reference system}
\label{sec:implications}

The two design laws translate directly into specification guidance, stated as \Modelled{}
guidance rather than a certified budget.

\emph{On the frame/offset separation.} A LunaNet or Moonlight/LCNS realization
\citep{israel2020lunanet,fienga2024moonlight,lnis2023} that fields a geometrically diverse
beacon and orbiter network can separate the lunar frame scale from the timescale offset with
one-way ranging alone (Design law~1); it need not mandate two-way ranging \emph{for datum
separation} (two-way ranging remains valuable for other reasons---absolute range, clock
synchronization---outside this analysis). The requirement to add a two-way or external time
tie applies specifically to the low-diversity / single-user regime, where the classical
radial/clock degeneracy bites.

\emph{On the timescale rate.} The IAU lunar-reference-system and lunar-coordinate-time
effort, and the Moonlight lunar-time work \citep{fienga2024moonlight}, must fix the rate by
\emph{convention}: the reference potential $W_0$, the selenoid surface, and the averaging
window (Design law~2). This is not a shortfall of the PNT system---no observation can reduce
the definitional freedom---but a statement about what a standard is \emph{for}. The
consolation is that the frame estimation, however imperfect, perturbs the rate by
$\lesssim\!\sim\!2\times10^{-17}$, negligible against the definitional band and below
near-term deployed optical-clock stability; the rate is robust to estimation error and needs
only a defined convention.

\emph{On the joint realization.} Theorem~\ref{thm:system} is the compact requirement: a
lunar reference \emph{system} standard must fix both a spatial datum and a temporal
convention; fixing one and estimating the other leaves an unobservable or undefined gauge.
This is the coupled-gauge counterpart of the spatial design law of \citep{baweja2026datum}
and the interoperability design law of \citep{baweja2026interop}.

\subsection{Limitations}
\label{sec:limitations}

The firewall of Section~\ref{sec:anchor} is also a statement of limitations; we make them
explicit.

\begin{itemize}
  \item \textbf{Modelled geometry.} The beacon and orbiter geometry is a representative
  construction, not a flown constellation. The classification (Theorem~\ref{thm:classify})
  and its counts are algebraic and hold for any network of the stated form; the specific
  eigenvalues and marginal values (Tables~\ref{tab:coupling1}--\ref{tab:anchor}) depend on
  the modelled geometry and would shift for a real network.

  \item \textbf{First-order linearization.} The datum, the observation rows, and the Fisher
  are first-order about a nominal. Second-order geometric effects (e.g.\ the finite-rotation
  residual, small on the P3 nanoradian scale) and non-Gaussian error structure are outside
  the model.

  \item \textbf{Coupling magnitudes are order-of-magnitude engineering bounds.} The three
  Jacobian entries \eqref{eq:dads}--\eqref{eq:dadr} and the propagated bound
  \eqref{eq:framebound} are first-principles, cross-checked to the published band, and
  unit-tested for internal consistency---but they are \emph{not} certified budgets. The
  definitional-band decomposition (tidal, $J_2$, $W_0$/selenoid, window) uses
  physics-literature reference values for the named sub-terms.

  \item \textbf{The coupled branch is proven, not externally validated.} The single
  external row (Section~\ref{sec:anchor}) exercises defect-$0$ and exact-rate-defect
  networks, both with $\operatorname{coupled\,dim}=0$. The nonzero-coupling classification is
  proven by construction and unit-tested (including rotated-basis invariance), not confirmed
  against an external dataset.

  \item \textbf{No flight data, no certification, no TRL claim.} Nothing here is a certified
  timescale realization, a certified frame, or a statement above analysis. There is no
  single-number operational rate claim; the rate is characterized by a definitional band and
  a negligible estimation coupling, not a certified value.
\end{itemize}

None of these caveats touches the proven content---the classification, its basis invariance,
the exact defect count, the necessary-not-sufficient criterion, and the lift theorem---or
the validated content---the linear-algebra reproduction on real DE440-orientation rows. Each
marks where a real lunar deployment must supply flight data and a fixed convention before any
absolute number is trusted.
